\begin{circuitikz}[
    american,
    voltage dir=RP,
    >=latex,
    %>={Latex[length=5mm, width=2mm]},
    alt={Ladder logic diagram spanning two vertical power rails, featuring a main rung with two parallel execution blocks: a MOVE block and a CPT compute block. The MOVE block receives Source Local:1:I.Ch0Data with value 2957 and copies it to Dest RawAnalog with value 2957. A parallel branch leads to the CPT block, which computes the expression open parenthesis RawAnalog plus 112.3 close parenthesis times open parenthesis 1 divided by 42.277 close parenthesis, storing the output in Dest TempF with value 72.5998. Green highlighted logic tracks power flow along the left vertical rail, into the MOVE block input, and down the branch into the CPT block input.}
    ]

    % Define the number of rungs in the diagram
    \def\numRungs{2}

    % Calculate the length of the vertical power rails dynamically based on number of rungs
    \pgfmathsetmacro{\railLength}{-2 * \numRungs}

    %======================================================================================================
    % Component Grid
    % Spacing = 3x, 2y
    %======================================================================================================
    \foreach \i in {0,...,5} {
        \foreach \j in {0,...,20} { % Change to 0,...,19 if you need 20 individual Y points down to -40
        
        % Calculate spatial positions
        \pgfmathsetmacro{\px}{\i * 3}
        \pgfmathsetmacro{\py}{-1 - (\j * 2)}
        
        % Name coordinate using loop indices: (C-column-row) e.g., (C-0-0), (C-5-10)
        \coordinate (C-\i-\j) at (\px, \py);
        
        % OPTIONAL: Uncomment the 2 lines below to display grid labels while building
        % \fill[red] (C-\i-\j) circle (1.5pt);
        %\node[anchor=south west, scale=0.5, gray] at (C-\i-\j) {(\i,\j)};
        }
    }

    %======================================================================================================
    % Foreground Layer
    % Only contains the vertical power rails
    % Length of the rails is automatic, calculated based on the number of rungs defined above
    %======================================================================================================
    \begin{pgfonlayer}{ft}

        %======================================================================================================
        % Scope sets options for both lines - BakerBlack, very thick
        %======================================================================================================
        \begin{scope}[
            draw=BakerBlack,
            text=BakerBlack,
            color=BakerBlack,
            font=\small,
            ultra thick
        ]
        
            \draw (0,0) -- (0,\railLength);

            \draw (15,0) -- (15,\railLength);

        \end{scope}
    \end{pgfonlayer}

    %======================================================================================================
    % Main Layer
    % All actual logic goes here
    %======================================================================================================
    \begin{pgfonlayer}{main}
        
        %======================================================================================================
        % Scope sets options for both lines - BakerBlack for draw and text
        %======================================================================================================
        \begin{scope}[
            draw=BakerBlack,
            text=BakerBlack,
            color=BakerBlack,
            font=\small
        ]

            %=================================================================
            % Rung 1
            %=================================================================
            \draw (C-0-0) -- (C-2-0) pic (MOV1) {SrcSingle};
            \draw (MOV1-OUT) -- (C-5-0);

            \node[anchor=south] at (MOV1-TOP) {MOVE};
            \node[anchor=west] at (MOV1-SRC) {\texttt{Local:1:I.Ch0Data}};
            \node[anchor=west] at (MOV1-SRCV) {2957};
            \node[anchor=west] at (MOV1-DEST) {\texttt{RawAnalog}};
            \node[anchor=west] at (MOV1-DESTV) {2957};

            \draw ($(C-1-0)+(1,0)$) coordinate(h1) -- ++(0,-3) coordinate(h2) -- ++(0.5,0) pic (CPT) {CPT};
            \draw (CPT-OUT) -- ++(0.5,0) coordinate(h3) -- ++(0,3) coordinate(h4);

            \node[anchor=south] at (CPT-TOP) {CPT};
            \node[anchor=west] at (CPT-DEST) {\texttt{TempF}};
            \node[anchor=west] at (CPT-DESTV) {72.5998};
            \node[anchor=west] at (CPT-EXP) {(\texttt{RawAnalog} + 112.3) * (1 / 42.277)};

        \end{scope}
    \end{pgfonlayer}

    %=================================================================
    % Background layer
    %=================================================================
    \begin{pgfonlayer}{bg}
        
        %==================================================================
        % Highlight Scope
        % This layer exists for highlighting true logic, the scope sets the
        % highlight options without having to repeat them every \draw
        % draw=UpNorthGreen, line width=10pt, opacity=0.4
        %==================================================================
        \begin{scope}[
            draw=UpNorthGreen,
            line width=10pt,
            opacity=0.4
        ]

            %==================================================================
            % Left vertical rung - leave this alone, it should ALWAYS be green
            % It automatically draws itself to the length defined at the top
            %==================================================================
            \draw (0,0) -- (0,\railLength);
        
            %=================================================================
            % Rung 1
            %=================================================================
            \draw (C-0-0) -- (MOV1-IN);
            % \draw (MOV1-OUT) -- (C-5-0);

            \draw (h1) -- (h2) -- (CPT-IN);
            % \draw (CPT-OUT) -- (h3) -- (h4);

        \end{scope}
    \end{pgfonlayer}

\end{circuitikz}